\begin{oving}
\begin{itemize}[leftmargin=5.6em,labelwidth=5em,labelsep=0.6em,align=left]
\item[\lett] \href{https://github.com/gustavbv/Matte4C---eksamen/blob/main/oppgavesett/F12-minste-kvadrater/3-hele-oppgaven-i-pensum/009-TMA4115-2025-var-oppg2-oppgave.pdf}{TMA4115 vår 2025, oppg.~2} {\color{graa}(\href{https://github.com/gustavbv/Matte4C---eksamen/blob/main/oppgavesett/F12-minste-kvadrater/3-hele-oppgaven-i-pensum/009-TMA4115-2025-var-oppg2-svar.pdf}{LF})}: rett linje gjennom tre punkter
\item[\middels] \href{https://github.com/gustavbv/Matte4C---eksamen/blob/main/oppgavesett/F12-minste-kvadrater/1-TMA4432-ovinger-og-proveeksamen/002-TMA4432-oving-1-oppg3-oppgave.pdf}{Øving 1, oppg.~3} {\color{graa}(\href{https://github.com/gustavbv/Matte4C---eksamen/blob/main/oppgavesett/F12-minste-kvadrater/1-TMA4432-ovinger-og-proveeksamen/002-TMA4432-oving-1-oppg3-svar.pdf}{LF})}: lineær regresjon med Moore--Penrose
\item[\middels] \href{https://github.com/gustavbv/Matte4C---eksamen/blob/main/oppgavesett/F12-minste-kvadrater/3-hele-oppgaven-i-pensum/031-TMA4115-2018-kont-oppg9-oppgave.pdf}{TMA4115 kont.\ 2018, oppg.~9} {\color{graa}(\href{https://github.com/gustavbv/Matte4C---eksamen/blob/main/oppgavesett/F12-minste-kvadrater/3-hele-oppgaven-i-pensum/031-TMA4115-2018-kont-oppg9-svar.pdf}{LF})}: andregradspolynom tilpasset fem punkter
\item[\vanskelig] \href{https://github.com/gustavbv/Matte4C---eksamen/blob/main/oppgavesett/F12-minste-kvadrater/2-Bergmann-eksamener/006-TMA4215-2021-host-oppg1-oppgave.pdf}{TMA4215 høst 2021, oppg.~1} {\color{graa}(\href{https://github.com/gustavbv/Matte4C---eksamen/blob/main/oppgavesett/F12-minste-kvadrater/2-Bergmann-eksamener/006-TMA4215-2021-host-oppg1-svar.pdf}{LF})}: QR med Givens-rotasjoner og minste kvadrater
\item[\vanskelig] \href{https://github.com/gustavbv/Matte4C---eksamen/blob/main/oppgavesett/F12-minste-kvadrater/2-Bergmann-eksamener/005-TMA4215-2022-host-oppg2-oppgave.pdf}{TMA4215 høst 2022, oppg.~2} {\color{graa}(\href{https://github.com/gustavbv/Matte4C---eksamen/blob/main/oppgavesett/F12-minste-kvadrater/2-Bergmann-eksamener/005-TMA4215-2022-host-oppg2-svar.pdf}{LF})}: gradient descent sammenlignet med normallikningen
\end{itemize}
\end{oving}
